% 三视图点修: 仅横向并排(正/侧/俯左至右)+图名下置居中基线对齐; 投影边线/尺寸/虚实/比例均沿用既有稿
\begin{tikzpicture}[line width=.6pt,font=\normalsize,>=Latex]
% 正视图 (宽2高4)
\begin{scope}[shift={(0.6,1.2)}]
\draw (0,0) -- (2,0) -- (2,4) -- (0,2) -- cycle;
% 底部宽2
\draw[thin] (0,0) -- (0,-0.5);
\draw[thin] (2,0) -- (2,-0.5);
\draw[thin,<->] (0,-0.3) -- (2,-0.3);
\node[below] at (1,-0.3) {2};
% 右侧高4
\draw[thin] (2,0) -- (2.5,0);
\draw[thin] (2,4) -- (2.5,4);
\draw[thin,<->] (2.3,0) -- (2.3,4);
\node[right] at (2.3,2) {4};
% 左侧高2
\draw[thin] (0,0) -- (-0.5,0);
\draw[thin] (0,2) -- (-0.5,2);
\draw[thin,<->] (-0.3,0) -- (-0.3,2);
\node[left] at (-0.3,1) {2};
\end{scope}
\node at (1.6,0) {正视图};
% 侧视图: 下矩宽2高2,上圆直径2,总高4与正视平齐
\begin{scope}[shift={(3.6,1.2)}]
\draw (0,0) -- (0,3);
\draw (2,0) -- (2,3);
\draw (0,0) -- (2,0);
\draw (1,3) circle (1);
% 底部宽2
\draw[thin] (0,0) -- (0,-0.5);
\draw[thin] (2,0) -- (2,-0.5);
\draw[thin,<->] (0,-0.3) -- (2,-0.3);
\node[below] at (1,-0.3) {2};
\end{scope}
\node at (4.6,0) {侧视图};
% 俯视图: 直径2圆,与正视同宽
\begin{scope}[shift={(6.4,1.2)}]
\draw (1,1) circle (1);
\end{scope}
\node at (7.4,0) {俯视图};
\end{tikzpicture}
